% Paper II skeleton -- structure only; text is written once Gate B numbers are
% frozen (report section 4.20). Build: pdflatex + bibtex with ../paper/references.bib.
\documentclass[fleqn,usenatbib]{mnras}
\usepackage{graphicx,amsmath,booktabs}
\title[No scalar thermalisation reproduces lanthanide fluorescence]{Line identity is the closure: direct fluorescence against the two-level atom in a lanthanide line forest}
\author[Y. Zhu]{Yonglin Zhu}
\begin{document}
\maketitle

\begin{abstract}
%% [to write] Same-code comparison of five+ treatments of the La II forest at
%% day 1, 1000-3000 km/s, 6000 K continuum: pure absorption (Sobolev,
%% expansion), complete thermalisation, scalar TLA(eps), direct A-beta
%% branching, and a branching-aware Poisson closure. Headline numbers from
%% report 4.19.3 / 4.20.1: refill 0.183 -> 0.660; eps=1 -38%, eps=0 +34%;
%% eps_best 0.02-0.36 by band, optical band unreachable; expansion_branch +21%.
\end{abstract}

\section{Introduction}
%% Narrative 1-7 (research plan): (1) Paper I: two approximations, not one,
%% attenuation only; (2) what re-emission changes; (3) the TLA and SN Ia
%% iron-peak experience (Kasen et al. 2006, eps ~ 0.3); (4) lanthanide
%% networks are not iron-peak networks; (5) the question: does any scalar eps
%% reproduce direct fluorescence; (6) the design: one code, same atom, same
%% rays, treatments differ by one choice each; (7) answer in one sentence.

\section{Methods}
\subsection{Same-code design}            %% forest_mc.py; lockstep packets; what is shared
\subsection{Sobolev resonances}           %% delta resonances, tau_S, beta
\subsection{Expansion opacity}            %% continuous absorber, inverted cumulative E_b, bins
\subsection{Direct $A\beta$ branching}    %% re-absorption (1-beta), exit kernel A beta / sum A beta
\subsection{Two-level atom: $\epsilon$}   %% thermal vs coherent per event; kappa_exp B_nu emissivity (E0)
\subsection{Packet bookkeeping}           %% photons carrying h nu; E_inj = E_esc + E_core + E_abs + E_dep (E1)
\subsection{The La II atom}               %% GSI 472 levels / 17,743 lines; 949 opacity lines 1148-17,609 A at 3000 K
\subsection{Validation against Paper I and SEDONA}   %% 0.3475/0.4839; RE sub-bands; escape probability (F19)

\section{Validation}
%% Table: Paper I legs; SEDONA RE sub-bands at 100 and 10 km/s (E2); energy identity; A-beta kernel test.

\section{Results}
\subsection{The sign reversal}                    %% 4.19.3 table; eps=1 -38%, eps=0 +34%
\subsection{Thermal-width convergence}            %% E2 table, SEDONA rows + one MC row
\subsection{The $\epsilon$ sweep}                 %% 4.20.1 tables; Fig fig_p2_eps_sweep
\subsection{Spectral $\chi^2(\epsilon)$}          %% 4.20.2; Fig fig_p2_eps_chi2
\subsection{Redistribution matrices}              %% 4.20.3; Fig fig_p2_redistribution
\subsection{Cascade pathways}                     %% E7: 3300-4500 A pumps, 5d6p, 971 pathways
\subsection{A line-dependent closure}             %% E8: expansion_branch +21%; bin-width
\subsection{Ce II and mixtures}                   %% placeholder (E9+)

\section{Discussion}
%% Opacity-representation error vs redistribution-closure error, separated by
%% the 2x2 {Sobolev, expansion} x {branch, TLA}; why eps ~ 0.3 "works" for
%% iron-peak and appears here only as a compensation; what a closure must
%% supply (the event kernel); TARDIS macroatom, ARTIS, SUMO; limits (LTE
%% populations, one ion, first-order Doppler, frozen T).

\section{Conclusions}

\bibliographystyle{mnras}
\bibliography{../paper/references}
\end{document}
